\documentclass[12pt,a4paper]{article}

\usepackage[utf8]{inputenc}
\usepackage[T1]{fontenc}
\usepackage[spanish]{babel}
\usepackage{amsmath,amssymb}
\usepackage{geometry}
\usepackage{graphicx}
\usepackage{float}
\usepackage{hyperref}

\geometry{margin=2.5cm}

\title{Reporte del Módulo de Cálculo de Hurst, DFA e Higuchi}
\author{Nombre del estudiante}
\date{\today}

\begin{document}

\maketitle

\section{Introducción}

En esta práctica se implementó una ventana dentro de la aplicación web para analizar series de tiempo financieras mediante cuatro métricas:

\begin{itemize}
    \item exponente de Hurst con rango reescalado,
    \item exponente de Hurst por escalamiento de incrementos,
    \item exponente DFA,
    \item dimensión fractal de Higuchi.
\end{itemize}

La pestaña trabaja con las series del mes calendario anterior para los activos VOO, NVDA y GAP. Además, incluye validaciones sintéticas independientes con ruido blanco y movimiento browniano.

\section{Objetivo}

El objetivo del módulo es estimar dos versiones del exponente de Hurst y compararlas con DFA e Higuchi sobre series de tiempo reales y sintéticas.

En particular:

\begin{enumerate}
    \item se implementó el cálculo de Hurst con el método de rango reescalado \(R/S\),
    \item se implementó el cálculo de Hurst por escalamiento de incrementos,
    \item se implementó DFA,
    \item se implementó la dimensión fractal de Higuchi,
    \item se aplicó el análisis a datos financieros reales,
    \item se incorporaron validaciones con ruido blanco y movimiento browniano.
\end{enumerate}

\section{Fundamento teórico}

\subsection{Hurst con rango reescalado}

En la versión final de la pestaña, el primer estimador de Hurst se calcula con el método clásico de rango reescalado.

Sea una serie \(X_1, X_2, \dots, X_n\). Primero se calcula el promedio:

\[
\bar X_n = \frac{1}{n}\sum_{i=1}^{n} X_i
\]

Luego se construye el perfil acumulado:

\[
Y(k)=\sum_{i=1}^{k}\bigl(X_i-\bar X_n\bigr), \qquad 1\le k\le n
\]

El rango del perfil es:

\[
R(n)=\max_{1\le k\le n} Y(k)-\min_{1\le k\le n} Y(k)
\]

La desviación estándar se define como:

\[
S(n)=\sqrt{\frac{1}{n}\sum_{i=1}^{n}\bigl(X_i-\bar X_n\bigr)^2}
\]

La ley de escalamiento usada es:

\[
\frac{R(n)}{S(n)} \sim c\, n^H
\]

Tomando logaritmos:

\[
\log\!\left(\frac{R(n)}{S(n)}\right)=\log c + H \log n
\]

Por lo tanto, el valor de \(H\) se obtiene como la pendiente del plano log-log:

\[
\log\!\left(\frac{R(n)}{S(n)}\right)
\quad \text{vs} \quad
\log n
\]

\subsection{Hurst por escalamiento de incrementos}

El segundo estimador de Hurst corresponde al primer método que se había construido en la práctica. En este caso se toma la serie \(X_t\) y se define:

\[
S(\lambda)=\mathrm{std}(X_{t+\lambda}-X_t)
\]

\[
S(\lambda)\propto \lambda^H
\]

\[
\log S(\lambda)= H \log \lambda + C
\]

La pendiente en el plano log-log permite estimar el segundo valor de \(H\). Este estimador sí se calcula para todas las series y validaciones, pero no se grafica en la parte inferior de la ventana final.

\subsection{Detrended Fluctuation Analysis}

El DFA se implementó siguiendo los pasos estándar.

Primero se integra la serie:

\[
Y(i)=\sum_{j=1}^{i}(X_j-\bar X)
\]

Luego se divide el perfil en ventanas de tamaño \(s\), sin traslape:

\[
N_s = \frac{N}{s}
\]

En cada ventana se ajusta una tendencia local lineal y se calcula:

\[
F_s^2=\frac{1}{s}\sum_{i=1}^{s}\bigl(Y_i-Y_i^{fit}\bigr)^2
\]

Posteriormente:

\[
F(s)=\sqrt{\frac{1}{N_s}\sum F_s^2}
\]

La ley de escalamiento es:

\[
F(s)\sim s^\alpha
\]

y por tanto:

\[
\log F(s)=\alpha \log s + C
\]

La pendiente de la recta en el plano log-log corresponde al exponente DFA.

\subsection{Dimensión fractal de Higuchi}

Para una serie \(X_t\), se define:

\[
X_k^m=\{X(m),X(m+k),X(m+2k),\dots,X(m+\lfloor (N-m)/k \rfloor k)\}
\]

La longitud asociada a cada \(m\) y \(k\) es:

\[
L_m(k)=
\frac{N-1}{\left\lfloor \frac{N-m}{k}\right\rfloor k^2}
\sum_{i=1}^{\left\lfloor \frac{N-m}{k}\right\rfloor}
\left|X(m+ik)-X(m+(i-1)k)\right|
\]

Promediando sobre \(m\):

\[
L(k)=\frac{1}{k}\sum_{m=1}^{k}L_m(k)
\]

La ley de escalamiento es:

\[
L(k)\propto k^{-D}
\]

y entonces:

\[
\log L(k)= -D\,\log k + C
\]

La dimensión de Higuchi se obtiene como el negativo de la pendiente.

\section{Descripción del código implementado}

El módulo fue implementado principalmente en el archivo \texttt{hurst.py} y su interfaz visual en \texttt{templates/hurst.html}. El flujo general del código es el siguiente.

\subsection{1. Selección automática del periodo}

El programa calcula automáticamente el mes calendario anterior a la fecha actual. Esto permite trabajar siempre con un periodo reciente sin necesidad de introducir fechas manualmente.

\subsection{2. Descarga de datos en tiempo real}

Las series financieras se obtienen desde Yahoo Finance mediante su endpoint de tipo \texttt{chart}. En esta implementación se usan:

\begin{itemize}
    \item VOO,
    \item NVDA,
    \item GAP.
\end{itemize}

\subsection{3. Estimación de Hurst}

Para cada serie se calculan dos exponentes de Hurst:

\begin{enumerate}
    \item Hurst con rango reescalado \(R/S\),
    \item Hurst por escalamiento de incrementos.
\end{enumerate}

El primero es el que se usa en la gráfica log-log inferior. El segundo se reporta numéricamente en las tarjetas resumen y en la tabla de resultados.

\subsection{4. Validaciones sintéticas}

El sistema genera dos validaciones:

\begin{itemize}
    \item ruido blanco con \texttt{np.random.normal(0,1)},
    \item movimiento browniano con \texttt{cumsum(np.random.normal(0,1))}.
\end{itemize}

Estas dos series se usan para validar ambos estimadores de Hurst, además de DFA e Higuchi.

\subsection{5. Cálculo de DFA e Higuchi}

Sobre las mismas series también se calculan:

\begin{itemize}
    \item el exponente DFA,
    \item la dimensión fractal de Higuchi.
\end{itemize}

\subsection{6. Visualización en la página web}

La ventana de Hurst hace lo siguiente:

\begin{itemize}
    \item muestra el periodo analizado,
    \item grafica las series financieras del mes anterior,
    \item grafica en ventanas independientes el ruido blanco y el movimiento browniano,
    \item muestra abajo solo el plano log-log de Hurst \(R/S\),
    \item muestra abajo el plano log-log de DFA,
    \item muestra abajo el plano log-log de Higuchi,
    \item construye una tabla con ambos Hurst, DFA e Higuchi.
\end{itemize}

\section{Funciones principales del módulo}

A nivel lógico, el código se organiza en estas partes:

\begin{enumerate}
    \item \texttt{previous\_calendar\_month}: determina el mes anterior automáticamente.
    \item \texttt{fetch\_yahoo\_daily\_series}: descarga los cierres diarios desde Yahoo Finance.
    \item \texttt{estimate\_hurst\_rs}: calcula \(H\) usando rango reescalado.
    \item \texttt{estimate\_hurst\_increment\_scaling}: calcula \(H\) por incrementos.
    \item \texttt{estimate\_dfa}: calcula el exponente \(\alpha\).
    \item \texttt{estimate\_higuchi\_dimension}: calcula la dimensión \(D\).
    \item \texttt{build\_white\_noise\_validation}: genera la validación con ruido blanco.
    \item \texttt{build\_brownian\_validation}: genera la validación con movimiento browniano.
    \item \texttt{build\_market\_payload}: reúne los datos de mercado y de validación.
    \item \texttt{register\_hurst}: registra las rutas \texttt{/hurst} y \texttt{/process\_hurst}.
\end{enumerate}

\section{Qué hace la ventana de Hurst}

La ventana nueva de la aplicación cumple con los siguientes objetivos:

\begin{itemize}
    \item cargar automáticamente información reciente del mercado,
    \item mostrar visualmente la evolución temporal de VOO, NVDA y GAP,
    \item calcular dos exponentes de Hurst para cada serie,
    \item comparar los resultados con ruido blanco y movimiento browniano,
    \item mostrar solo el Hurst \(R/S\) en la gráfica inferior,
    \item dejar el Hurst por incrementos como valor numérico no graficado.
\end{itemize}

\section{Resultados esperados}

Si la implementación es correcta, entonces:

\begin{itemize}
    \item con \(R/S\), el ruido blanco debe acercarse a \(H \approx 0.5\),
    \item con \(R/S\), el movimiento browniano debe acercarse a \(H \approx 1\),
    \item con el método por incrementos, el ruido blanco debe acercarse a \(H \approx 0\),
    \item con el método por incrementos, el movimiento browniano debe acercarse a \(H \approx 0.5\),
    \item para DFA, el ruido blanco debe acercarse a \(\alpha \approx 0.5\),
    \item para DFA, el movimiento browniano debe acercarse a \(\alpha \approx 1.5\),
    \item para Higuchi, el ruido blanco debe acercarse a \(D \approx 2\),
    \item para Higuchi, el movimiento browniano debe acercarse a \(D \approx 1.5\).
\end{itemize}

En muestras finitas pueden aparecer pequeñas desviaciones, pero la tendencia general debe respetar esas referencias.

\section{Espacios para capturas}

\subsection{Ventana completa del módulo}

\begin{figure}[H]
    \centering
    \fbox{\rule{0pt}{7cm}\rule{0.9\linewidth}{0pt}}
    \caption{Vista general de la pestaña de Hurst, DFA e Higuchi.}
\end{figure}

\subsection{Series financieras del mes anterior}

\begin{figure}[H]
    \centering
    \fbox{\rule{0pt}{7cm}\rule{0.9\linewidth}{0pt}}
    \caption{Series financieras de VOO, NVDA y GAP.}
\end{figure}

\subsection{Ruido blanco}

\begin{figure}[H]
    \centering
    \fbox{\rule{0pt}{7cm}\rule{0.9\linewidth}{0pt}}
    \caption{Ventana independiente para la validación con ruido blanco.}
\end{figure}

\subsection{Movimiento browniano}

\begin{figure}[H]
    \centering
    \fbox{\rule{0pt}{7cm}\rule{0.9\linewidth}{0pt}}
    \caption{Ventana independiente para la validación con movimiento browniano.}
\end{figure}

\subsection{Plano log-log de Hurst \(R/S\)}

\begin{figure}[H]
    \centering
    \fbox{\rule{0pt}{7cm}\rule{0.9\linewidth}{0pt}}
    \caption{Plano log-log de \(\log(R/S)\) contra \(\log(n)\), donde la pendiente es \(H\).}
\end{figure}

\subsection{Plano log-log de DFA}

\begin{figure}[H]
    \centering
    \fbox{\rule{0pt}{7cm}\rule{0.9\linewidth}{0pt}}
    \caption{Plano log-log para DFA sobre ruido blanco.}
\end{figure}

\subsection{Plano log-log de Higuchi}

\begin{figure}[H]
    \centering
    \fbox{\rule{0pt}{7cm}\rule{0.9\linewidth}{0pt}}
    \caption{Plano log-log para la dimensión de Higuchi sobre ruido blanco.}
\end{figure}

\subsection{Tabla final}

\begin{figure}[H]
    \centering
    \fbox{\rule{0pt}{6cm}\rule{0.9\linewidth}{0pt}}
    \caption{Tabla con Hurst \(R/S\), Hurst por incrementos, DFA e Higuchi para mercado y validaciones sintéticas.}
\end{figure}

\section{Conclusión}

La ventana desarrollada permite estudiar memoria, persistencia y rugosidad de series de tiempo reales y sintéticas desde un mismo entorno web. El módulo integra descarga automática de datos, cálculo de dos exponentes de Hurst, validaciones con ruido blanco y browniano, y visualización en gráficas y tablas.

En particular, el sistema permite corroborar de forma directa que:

\begin{itemize}
    \item el método \(R/S\) estima \(H\) como la pendiente del plano \(\log(R/S)\) contra \(\log(n)\),
    \item el método por incrementos recupera la validación con ruido blanco cercana a \(H \approx 0\),
    \item DFA distingue entre ruido blanco y browniano,
    \item Higuchi distingue la rugosidad entre ambos procesos.
\end{itemize}

\end{document}
